\section{The primitive minimal class}
\label{sec:minimal-class}

Let \((A,\Theta)\) be a principally polarized abelian variety of dimension
\(g\).  Write
\[
 \gamma_A=\frac{\Theta^{g-1}}{(g-1)!}\in H^{2g-2}(A,\Z).
\]
The issue is not the rational algebraicity of \(\gamma_A\), but its membership
in the integral lattice generated by ordinary products of divisor classes.

\begin{theorem}[Semisimple graph slopes]
\label{thm:semisimple-slope}
Let \(p\) be a prime and let \(T\in\operatorname{Sym}_g(\F_p)\).  Form the
principally polarized quotient of the elliptic-power source with coefficient
polarization \(pI_g\) by the maximal isotropic graph of \(T\).  If the
minimal polynomial of \(T\) is squarefree, then
\[
 \gamma_A\in
 \operatorname{im}\bigl(\operatorname{Sym}^{g-1}\NS(A)
       \longrightarrow H^{2g-2}(A,\Z)\bigr).
\]
\end{theorem}

\begin{proof}
Work over \(\Z_p\), and pass to a finite unramified extension over which the
squarefree minimal polynomial splits.  Self-adjoint eigenspaces for distinct
roots are orthogonal and nondegenerate.  They lift, without dividing by two,
to an orthogonal decomposition with blocks \(pB_\lambda\), where each
\(B_\lambda\) is symmetric and unimodular.

On a scalar slope block, the integral divisor lattice contains
\(p\operatorname{Sym}(B_\lambda)\).  Fill every non-target block by its
primitive mixed determinant.  On the target block, the degree-one-lower
mixed adjugate gives every diagonal matrix unit and every symmetric
off-diagonal matrix unit with coefficient \(\pm1\).  Summing over target
blocks constructs the cofactor of
\(\mathop{\perp}_{\lambda}pB_\lambda\), which is
the pullback of \(\gamma_A\).  No factorial or off-diagonal factor two is
introduced.  Membership descends from the unramified extension by faithful
flatness.  Away from \(p\), the quotient isogeny identifies the integral
divisor and cohomology lattices, so local membership globalizes.
\end{proof}

The scalar version of the same mixed-adjugate construction applies to the
three-primary block of the six-axis polarization.  We can therefore use the
two local forms isolated in Section~\ref{sec:envelope}.

\begin{lemma}[Local chart for the six-axis quotient]
\label{lem:six-axis-local-chart}
Let \(G=6I_5-J_5\).  At each prime \(p\mid6\), the coefficient lattice has an
orthogonal Jordan decomposition
\[
  (\Z_p^5,G)\simeq U_0\perp pU_1,
\]
where \(U_0\) has rank one and is unimodular and \(U_1\) has rank four and
is unimodular.  The principal quotient is trivial on \(U_0\).  On \(U_1/pU_1\),
its graph slope has minimal polynomial \(t^2+t+1\) at \(p=2\), while at
\(p=3\) it is scalar on the simple coefficient constituent.
\end{lemma}

\begin{proof}
The Smith form of \(G\) is \((1,6,6,6,6)\).  The unit elementary divisor
splits off orthogonally over \(\Z_p\), leaving \(p\) times a unimodular
rank-four form.  The unit block has no discriminant and hence no gluing
choice.  At two the discriminant heart is the six-point module \(H_2\), and
Proposition~\ref{prop:golden-exotic} identifies the selected endomorphism
with \(\omega\in\F_4\); its polynomial is \(t^2+t+1\).  At three the
coefficient heart is simple with scalar commutant, and the unique
monodromy-invariant elliptic three-torsion line makes the graph scalar.
\end{proof}

\begin{theorem}[Minimal class for the exotic cubic family]
\label{thm:exotic-minimal-class}
For every smooth member \(X_b\) of the nonstandard \(A_5\)-cubic pencil, with
intermediate Jacobian \((J_b,\Theta_b)\),
\[
 \frac{\Theta_b^4}{4!}
 \in
 \operatorname{im}\bigl(\operatorname{Sym}^4\NS(J_b)
       \longrightarrow H^8(J_b,\Z)\bigr).
\]
\end{theorem}

\begin{proof}
The five-axis source has polarization type \((1,6,6,6,6)\), so only the
primes two and three can contribute an integral defect.  At two the exotic
graph slope acts on the simple heart through \(\omega\in\F_4\); its minimal
polynomial \(t^2+t+1\) is squarefree.  Theorem~\ref{thm:semisimple-slope}
therefore gives the primitive local class.  At three the monodromy line makes
the gluing scalar on the simple coefficient block, and the scalar
mixed-adjugate formula gives the same conclusion.  Away from two and three
the source and quotient lattices are unimodular.  Local membership at every
prime is equivalent to integral membership in the divisor-product lattice.

Lemma~\ref{lem:six-axis-local-chart} identifies the actual geometric lattice,
not merely an abstract gluing with the same Smith form.  The graph kernel is a
finite sub-local-system on each connected smooth component of the pencil, so
the same local presentation holds on every fibre.  At a fibre with additional
N\'eron--Severi classes the divisor-product image can only increase.
\end{proof}

Voisin identifies this primitive class as the exact obstruction relevant to
cubic threefolds.

\begin{corollary}[Universal \(CH_0\)-triviality]
\label{cor:universal-ch0}
Every smooth member of the nonstandard \(A_5\)-cubic pencil is universally
\(CH_0\)-trivial.
\end{corollary}

\begin{proof}
By Theorem~\ref{thm:exotic-minimal-class}, the primitive minimal class of the
intermediate Jacobian is algebraic.  Voisin's Corollary~4.4
\cite{Voisin} says that, for a smooth complex cubic threefold, this is
equivalent to universal \(CH_0\)-triviality.
\end{proof}

The argument is fibrewise.  It constructs neither a horizontal minimal cycle
over the pencil nor a relative decomposition of the diagonal.  Neither is
needed for Corollary~\ref{cor:universal-ch0}.
